\TieudeT{PHONG TỎA VẬT LÍ 12 - VẬT LÍ NHIỆT}
\subsubsection*{PHẦN I. Thí sinh trả lời câu hỏi từ câu 1 đến câu 18. Mỗi câu hỏi thí sinh chỉ chọn một phương án}
\setcounter{ex}{0}
\begin{ex} %Câu 1
	Quần áo ướt được phơi khô, đó là hiện tượng
	\choice
	{ngưng tụ}
	{\True bay hơi}
	{đông đặc}
	{nóng chảy}
	\loigiai{}
\end{ex}

\begin{ex} %Câu 2
	Hệ thức $\Delta U = A + Q$ với $A > 0$, $Q < 0$ diễn tả cho quá trình nào của chất khí?
	\choice
	{Nhận công và nội năng giảm}
	{Nhận nhiệt và sinh công}
	{Tỏa nhiệt và nội năng giảm}
	{\True Nhận công và tỏa nhiệt}
	\loigiai{}
\end{ex}

\begin{ex} %Câu 3
	Trường hợp nào dưới đây làm biến đổi nội năng của đồng xu không do truyền nhiệt?
	\choice
	{Thả vào nước sôi}
	{Hơ nóng trên bếp lửa}
	{Phơi đồng xu ngoài nắng}
	{\True Cọ xát đồng xu trên mặt bàn}
	\loigiai{}
\end{ex}

\begin{ex} %Câu 4
	Một chai nước suối ở thể lỏng sau đó cho chai nước vào tủ đông, sau một thời gian lấy chai nước ra nước trong chai bị đông cứng lại. Đây gọi là hiện tượng gì?
	\choice
	{Sự nóng chảy}
	{\True Sự đông đặc}
	{Sự ngưng tụ}
	{Sự bay hơi}
	\loigiai{}
\end{ex}

\begin{ex} %Câu 5
	Dùng nhiệt kế rượu (có giới hạn đo từ $-20^\circ C$ đến $50^\circ C$) không thể đo được nhiệt độ nào sau đây?
	\choice
	{\True Nhiệt độ của nước đang sôi}
	{Nhiệt độ của nước đá đang tan}
	{Nhiệt độ của không khí trong phòng}
	{Nhiệt độ của nước uống}
	\loigiai{}
\end{ex}

\begin{ex} %Câu 6
	Trong mạng tinh thể NaCl, các ion $Na^+$ và $Cl^-$ được phân bố đều đặn trên các đỉnh và tâm các mặt của
	\choice
	{\True hình lập phương}
	{hình tứ diện đều}
	{hình chóp tam giác}
	{hình lăng trụ tam giác đều}
	\loigiai{}
\end{ex}

\begin{ex} %Câu 7
	Nội năng của một vật gồm
	\choice
	{tổng động năng của các phân tử cấu tạo nên vật}
	{tổng thế năng của các phân tử cấu tạo nên vật}
	{\True tổng động năng và thế năng của các phân tử cấu tạo nên vật}
	{tổng động năng và thế năng của vật}
	\loigiai{}
\end{ex}

\begin{ex} %Câu 8
	Phát biểu nào sau đây là đúng?
	\choice
	{Nước được cấu tạo từ các nguyên tử nước}
	{Nguyên tử gồm một nhóm các phân tử kết hợp lại}
	{\True Sắt được cấu tạo từ các nguyên tử sắt}
	{Gỗ được cấu tạo từ nhiều nguyên tử gỗ}
	\loigiai{}
\end{ex}

\begin{ex} %Câu 9
	Khi khoảng cách giữa các phân tử rất nhỏ, thì giữa các phân tử
	\choice
	{chỉ có lực hút}
	{\True có cả lực hút và lực đẩy, nhưng lực đẩy lớn hơn lực hút}
	{chỉ có lực đẩy}
	{có cả lực hút và lực đẩy, nhưng lực đẩy nhỏ hơn lực hút}
	\loigiai{}
\end{ex}

\begin{ex} %Câu 10
	Phát biểu nào sau đây là \textbf{sai} khi nói về các cách làm thay đổi nội năng của một vật?
	\choice
	{Nội năng của vật có thể biến đổi bằng hai cách: thực hiện công và truyền nhiệt}
	{Quá trình làm thay đổi nội năng có liên quan đến sự chuyển dời của các vật khác tác dụng lực lên vật đang xét gọi là sự thực hiện công}
	{Quá trình làm thay đổi nội năng không bằng cách thực hiện công gọi là sự truyền nhiệt}
	{\True Quá trình làm thay đổi nội năng bằng cách thực hiện công gọi là sự truyền nhiệt}
	\loigiai{}
\end{ex}


\begin{ex}%Câu 11
	Thả đồng thời thỏi đồng khối lượng $400\ g$, ở nhiệt độ $150^{\circ}C$ và khối sắt khối lượng $500\ g$, nhiệt độ $80^{\circ}C$ vào $2\ l$ nước ở nhiệt độ $25^{\circ}C$. Biết nhiệt dung riêng của đồng, sắt và nước lần lượt là $380\ J/kg.K$, $460\ J/kg.K$, $4200\ J/kg.K$ và coi như chỉ có 3 chất trao đổi nhiệt với nhau. Biết khối lượng riêng của nước bằng $1\ kg/dm^3$. Nhiệt độ cuối cùng của 3 chất là
	\choice
	{\True $28,6^{\circ}C$}
	{$38,6^{\circ}C$}
	{$48,6^{\circ}C$}
	{$58,6^{\circ}C$}
	\loigiai{
	\begin{itemize}
		\item Phương trình cân bằng nhiệt: $m_{Cu}.c_{Cu}(t_{Cu} - t_{cb}) + m_{Fe}.c_{Fe}(t_{Fe} - t_{cb}) + m_{n}.c_{n}(t_{n} - t_{cb}) = 0 \Rightarrow 0,4.380.(150-t_{cb}) + 0,5.460 (80- t_{cb}) + 2.4200(25-t_{cb}) =0\Rightarrow t_{cb} \approx 28,6^\circ C$.
	\end{itemize}
	}
\end{ex}

\begin{ex}%Câu 12
	Người ta đổ một lượng nước lạnh ở nhiệt độ $20^{\circ}C$ và một lượng nước nóng ở nhiệt độ $70^{\circ}C$ vào một bình cách nhiệt. Đến khi cân bằng nhiệt thì thu được nước ở nhiệt độ $30^{\circ}C$. Bỏ qua sự trao đổi nhiệt giữa nước với bình và với môi trường ngoài. Tỉ số giữa thể tích của nước lạnh và thể tích của nước nóng lúc ban đầu là
	\choice
	{\True $4$}
	{$5$}
	{$2,5$}
	{$3,5$}
	\loigiai{
	\begin{itemize}
		\item $m_l.c.(t-t_l) = m_n.c.(t_n-t) \Rightarrow V_l.(t-t_l) = V_n(t_n-t) \Rightarrow \dfrac{V_l}{V_n} = \dfrac{t_n-t}{t-t_n} = 4$.
	\end{itemize}
	}
\end{ex}

\begin{ex}%Câu 13
	\immini[thm]{Sự phụ thuộc của nhiệt độ vào nhiệt lượng tỏa ra của $0,2\ kg$ chất khí ban đầu được thể hiện trên hình. Quá trình này diễn ra ở áp suất không đổi. Nhiệt hóa hơi riêng của chất này là
		\choice
		{\True $3.10^4\ J/kg$}
		{$4.10^5\ J/kg$}
		{$4.10^4\ J/kg$}
		{$3,5.10^5\ J/kg$}}
		{
		\begin{tikzpicture}[draw=Mapcolor, very thick, >=latex']
			\draw [->] (0,0) -- (0,3.5) node [right] {$T$};
			\draw [->] (0,0) -- (4.5,0) node [above right] {$Q(kJ)$};
			\draw [dashed] (0,0) grid[xstep = 0.8, ystep = 0] (4,3);
			\draw (0,3) -- (0.8,1) -- (3.2,1) -- (4,0.5)
			node [left] at (0,0) {$O$};
			\foreach \x in {2,4,6,8}{
				\draw node [below] at (\x*0.4,0) {$\x$};}
		\end{tikzpicture}
		}
	
	\loigiai{
	\begin{itemize}
		\item $Q = m.L \Rightarrow L = \dfrac{Q}{m} = \dfrac{(8-2).10^3}{0,2} = 3.10^4\ J/kg$.
	\end{itemize}
	}
\end{ex}

\begin{ex}%Câu 14
	Một viên đạn chì va chạm vào vật cản cứng . Cho rằng 80\% động năng của viên đạn chuyển thành nội năng của nó khi va chạm; nhiệt độ của viên đạn trước khi va chạm là $127^\circ C$. Cho biết nhiệt dung riêng của chì là $c = 130\ J/kg.K$; nhiệt độ nóng chảy của chì là $327^\circ C$, nhiệt nóng chảy riêng của chì là $\lambda = 25\ kJ/kg$. Để khi va chạm vào vật cản cứng thì nóng chảy hoàn toàn thì viên đạn phải có tốc độ tối thiểu là
	\choice
	{\True $357\ m/s$}
	{$324\ m/s$}
	{$352\ m/s$}
	{$457\ m/s$}
	\loigiai{
		\begin{itemize}
			\item  Ta có: $Q = mc\Delta t + \lambda m \Rightarrow H.\dfrac{1}{2}mv^2 = mc\Delta t + \lambda m \Rightarrow 0,8.0,5v^2 = 130.(327-127) + 25.10^3 \Rightarrow v \approx 357\ m/s$.
		\end{itemize}
	}
\end{ex}

\begin{ex}%Câu 15
	Một học sinh cho rằng: "Nội năng của vật chỉ phụ thuộc vào nhiệt độ, vì khi nhiệt độ tăng thì các phân tử chuyển động nhanh hơn, làm tăng động năng của chúng." Cách lập luận này có đầy đủ chưa? Vì sao?
	\choice
	{Đúng, vì nội năng chỉ liên quan đến chuyển động của các phân tử trong vật}
	{\True Sai, vì nội năng không chỉ phụ thuộc vào nhiệt độ mà còn phụ thuộc vào thể tích do ảnh hưởng của thế năng phân tử}
	{Đúng, vì khi nhiệt độ không đổi thì nội năng không thể thay đổi}
	{Sai, vì nội năng chỉ phụ thuộc vào thể tích mà không liên quan đến nhiệt độ}
	\loigiai{
		\begin{itemize}
			\item Ta có: Nội năng $=$ Động năng của các phân tử $+$ thế năng phân tử.
			\item Mà động năng thì phụ thuộc nhiệt độ ($t$ tăng $\Leftrightarrow v$ tăng $\Leftrightarrow W_{\text{đ}}$ tăng...);
			\item Còn thế năng phân tử phụ thuộc thể tích (V thay đổi $\Rightarrow$ khoảng cách phân tử thay đổi $\Rightarrow$ thế năng tương tác phân tử thay đổi).
		\end{itemize}
		
	}
\end{ex}

\begin{ex}%Câu 16
 Một bạn học sinh làm thí nghiệm với đầy đủ thiết bị để xác định được nhiệt nóng chảy riêng của một chất khi đã biết nhiệt dung riêng của chất đó trong trạng thái rắn và trạng thái lỏng. Phát biểu nào sau đây là \textbf{sai}?
	\choice
	{Bắt đầu thí nghiệm từ khi chất đó đang ở trạng thái rắn và kết thúc khi chất đó đã hoàn toàn ở trạng thái lỏng}
	{\True Bắt đầu thí nghiệm từ khi chất đó đang ở trạng thái rắn và kết thúc khi đã thấy đã có sự nóng chảy của chất đó}
	{Thực hiện thí nghiệm từ khi chất bắt đầu đạt đến nhiệt độ nóng chảy nhưng chưa nóng chảy và kết thúc đo khi chất đó đã hoàn toàn ở trạng thái lỏng}
	{Thực hiện thí nghiệm từ khi chất đó bắt đầu đạt đến nhiệt độ nóng chảy nhưng chưa nóng chảy và kết thúc khi nóng chảy hoàn toàn và vẫn đang ở nhiệt độ nóng chảy}
	\loigiai{
		\begin{itemize}
			\item Chỉ bắt đầu nóng chảy chứ chưa xảy ra hiện tượng nóng chảy.
		\end{itemize}
	}
\end{ex}

\begin{ex}%Câu 17
	Dùng một ca múc nước ở thùng chứa nước $A$ có nhiệt độ $t_A = 20^\circ C$ và ở thùng chứa nước $B$ ở nhiệt độ $t_B = 80^\circ C$ rồi đổ vào thùng nước $C$. Biết rằng trước khi đổ, trong thùng chứa nước $C$ đã có sẵn một lượng nước ở nhiệt độ $t_C = 40^\circ C$ và bằng tổng số ca nước vừa mới đổ thêm vào nó. Bỏ qua sự trao đổi nhiệt với môi trường, với bình chứa và ca múc nước. Để có nhiệt độ nước ở thùng $C$ là $50^\circ C$ thì số ca nước phải múc ở mỗi thùng $A$ và $B$ lần lượt là
	\choice
	{$3 n$ ca và $2 n$ ca}
	{\True $n$ ca và $2 n$ ca}
	{$2 n$ ca và $n$ ca}
	{$n$ ca và $3 n$ ca}
	\loigiai{
		\begin{itemize}
			\item Phương trình cân bằng nhiệt: $Q_B = Q_A + Q_C \Rightarrow m_B.c(80-50) = m_A.c(50-20) + (m_A+m_B).c.(50-40) \Rightarrow m_B = 2m_A$.
		\end{itemize}
	}
\end{ex}
\begin{ex}%Câu 18
	Để diệt trừ các bào tử nấm và kích thích quá trình nảy mầm của hạt giống lúa, người nông dân đã sử dụng một kinh nghiệm dân gian là ngâm chúng vào trong nước ấm theo công thức "hai sôi, ba lạnh". Tức là nước ấm ở nhiệt độ $T_A$ ($^\circ C$) sẽ được tạo ra bằng cách pha hai phần nước sôi ($100^\circ C$) với ba phần nước lạnh $T_L$ ($^\circ C$). Phát biểu nào sau đây là đúng?
	\choice
	{ $T_A = \dfrac{200+3T_L}{2}$ }
	{ $T_A = \dfrac{200+2T_L}{3}$ }
	{\True $T_A = \dfrac{200+3T_L}{5}$ }
	{ $T_A = \dfrac{100+3T_L}{2}$ }
	\loigiai{
	\begin{itemize}
		\item Phương trình cân bằng nhiệt: $m_s.c.(T_s-T_A) = m_L.c.(T_A - T_L) \Rightarrow T_A = \dfrac{m_sT_s + m_LT_L}{m_s+m_L} = \dfrac{200+3T_L}{5}$.
	\end{itemize}
	}
\end{ex}
\subsubsection*{PHẦN 2. Thí sinh trả lời từ câu 1 đến câu 4. Trong mỗi ý a), b), c), d) ở mỗi câu, thí sinh chọn đúng hoặc sai.}
\setcounter{ex}{0}
\begin{ex}
	Khi nói về sự sôi và sự bay hơi.
	\choiceTF[t]
	{Sự sôi chỉ xảy ra trong lòng chất lỏng; sự bay hơi chỉ xảy ra trên bề mặt chất lỏng}
	{Sự sôi xảy ra cả bên trong và trên bề mặt chất lỏng; sự bay hơi chỉ xảy ra trong lòng chất lỏng}
	{\True Sự sôi xảy ra cả bên trong và trên bề mặt chất lỏng; sự bay hơi chỉ xảy ra trên bề mặt chất lỏng}
	{\True Sự sôi xảy ra ở nhiệt độ sôi của chất lỏng; sự bay hơi xảy ra ở mọi nhiệt độ}
	\loigiai{
		\begin{itemchoice}
			\itemch Sự sôi xảy ra ở cả bề mặt lẫn bên trong chất lỏng.
			\itemch Sự bay hơi chỉ xảy ra ở bề mặt của chất lỏng, không phải trong lòng chất lỏng.
			\itemch Sự sôi xảy ra ở cả bên trong và trên bề mặt chất lỏng; sự bay hơi chỉ xảy ra trên bề mặt chất lỏng.
			\itemch Sự sôi xảy ra ở nhiệt độ sôi của chất lỏng; sự bay hơi xảy ra ở mọi nhiệt độ.
		\end{itemchoice}
	}
\end{ex}

\begin{ex}
	\immini[thm]
	{Trong thí nghiệm đo nhiệt nóng chảy riêng của nước đá, người ta sử dụng dòng điện để cung cấp nhiệt lượng cho $0,6\ kg$ nước đá trong nhiệt lượng kế. Dùng Oát kế đo được công suất dòng điện cung cấp là $930\ W$, dùng đồng hồ để đo thời gian và nhiệt kế để đo nhiệt độ của khối nước. Đồ thị thực nghiệm đo được như hình.}{\begin{tikzpicture}[draw=Mapcolor, very thick, >=latex']
			\draw [dotted] (0,0) grid [xstep = 0.3, ystep = 0] (6,3);
			\draw [->] (0,0) -- (6.5,0) node [above] {$\tau (s)$};
			\draw [->] (0,0) -- (0,3) node [right] {$t (^\circ C)$};
			\draw [line width=1.5pt] (0,0) -- (3.3,0) -- (5.7,2.5);
			\draw
			(1.5,0) circle (1.5pt) node [below] {$100$}
			(3,0) circle (1.5pt) node [below]{$200$}
			(4.5,0) circle (1.5pt);
			\draw [dashed] (0,2.5) circle (1pt) node [left] {$100$} -- (5.7,2.5);
	\end{tikzpicture}}
	\choiceTFt
	{\True Thời gian để nước đá tan hoàn toàn là $220\ s$}
	{\True Nếu hao phí nhiệt lượng là $2\%$, thì nhiệt nóng chảy riêng của nước đá là $334,18\ kJ/kg$}
	{Bỏ qua mọi hao phí thì nhiệt nóng chảy riêng của nước đá là $310\ kJ/kg$}
	{Trong giai đoạn từ thời điểm $0\ s$ đến $200\ s$, khối chất không nhận thêm nhiệt lượng}
	\loigiai{
		\begin{itemchoice}
			\itemch Thời gian để nước đá tan hoàn toàn $200 + 1.\dfrac{100}{5} = 220\ s$.
			\itemch $Q = m.\lambda \Rightarrow \mathcal{H.P}.\tau = m.\lambda \Rightarrow \lambda = \dfrac{\mathbb{H.P}.t}{m} = \dfrac{(1-0,02).930.220}{0,6} = 334180\ J/kg = 334,18 kJ/kg$.
			\itemch $Q = m.\lambda \Rightarrow \mathcal{P}.\tau = m.\lambda \Rightarrow \lambda = \dfrac{\mathcal{P}.t}{m} = \dfrac{930.220}{0,6} = 341000\ J/kg = 341 kJ/kg$.
			\itemch Từ $0$ đến $220\ s$, nước đá nhận nhiệt lượng để nóng chảy.
		\end{itemchoice}
	}
\end{ex}

\begin{ex}
	Khi tiến hành nung nóng một chất rắn kết tinh bằng một bếp có công suất không đổi. Bỏ qua sự mất mát nhiệt lượng ra môi trường. Kể từ lúc bắt đầu đun người ta ghi nhận được đồ thị sự phụ thuộc của nhiệt độ của khối chất và thời gian đun như hình bên.
	\immini{\choiceTFt
	{Nhiệt độ nóng chảy của chất rắn trên là $327^\circ C$}
	{\True Kể từ lúc bắt đầu đun, nhiệt lượng cần để chất rắn tăng lên đến nhiệt độ nóng chảy gấp $1,56$ lần nhiệt lượng cần cung cấp trong suốt giai đoạn nóng chảy}
	{Đoạn AB trên đồ thị thể hiện quá trình chất rắn đang nóng chảy}
	{Tại phút thứ $39$ chất rắn đã nóng chảy hoàn toàn}}
	{\begin{tikzpicture}[draw=Mapcolor, very thick, >=latex']
			\draw [->] (0,0) -- (4.5,0) node [above] {$\tau$ (phút)};
			\draw [->] (0,0) -- (0,4) node [right] {$t(^\circ C)$};
			\draw (0,1) circle (1pt) node [below right] {$A$} -- (2,3) circle (1pt) node [above] {$B$} -- (3,3) circle (1pt) node [above] {$C$}-- (4,3.5) circle (1pt) node [above] {$D$};
			\draw [dashed] (0,3) circle (1pt) node [left] {$327$} -- (2,3) -- (2,0) circle (1pt) node [below] {$39$};
			\draw [dashed] (3,3) -- (3,0) node [below] {$64$}
			node [left] at (0,0) {$O$}
			node [left] at (0,1) {$27$};
	\end{tikzpicture}}
	\loigiai{
		\begin{itemchoice}
			\itemch Nhiệt độ nóng chảy không đổi là $327^\circ C$.
			\itemch $Q = \mathcal{P}.t \Rightarrow \dfrac{Q_{AB}}{Q_{BC}} = \dfrac{t_{AB}}{t_{BC}} = \dfrac{39}{64-39} = 1,56$.
			\itemch Trong đoạn $BC$ chất rắn đang nóng chảy.
			\itemch Tại phút thứ $39$ chất rắn mới đến điểm nóng chảy.
		\end{itemchoice}
	}
\end{ex}

\begin{ex}
	Giả sử một học sinh tạo ra một nhiệt kế sử dụng một thang nhiệt độ mới cho riêng mình, gọi là thang nhiệt độ $Z$, có đơn vị là $^\circ Z$. Trong đó, nhiệt độ của nước đá đang tan ở $1\ atm$ là $x^\circ Z$ và nhiệt độ nước sôi ở $1atm$ là $y^\circ Z$. Từ vạch $x^\circ Z$ đến vạch $y^\circ Z$ được chia thành $180$ khoảng, mỗi khoảng ứng với $1^\circ Z$.
	\choiceTFt
	{Một độ chia trên thang nhiệt độ $Z$ bằng $1,8$ lần độ chia trên thang nhiệt độ Celsius}
	{\True Mối liên hệ giữa $x$ và $y$ là: $y = x + 180$}
	{Độ biến thiên nhiệt độ $18^\circ$C trong thang nhiệt độ Celsius bằng với độ biến thiên nhiệt độ $10^\circ$Z trong thang nhiệt độ $Z$}
	{\True Nếu nhiệt độ cơ thể người là $37\ ^\circ C$ tương ứng với $86,6^\circ Z$ thì giá trị của $x$ là $20$}
	\loigiai{
		\begin{itemchoice}
			\itemch Một độ chia trên thang độ $Z$ bằng $\dfrac{100}{180}$ lần độ chia trên thang Celcius.
			\itemch $W = \dfrac{1}{2}mv_0^2 = 0,5.0,02.500^2 = 2500\ J$.
			\itemch Độ biến thiên $10^\circ C$ bằng $18^\circ Z$.
			\itemch $\dfrac{Z-x}{180} = \dfrac{t-0}{100-0} \Rightarrow \dfrac{86,6-x}{180} = \dfrac{37}{100-0} \Rightarrow x = 20$.
		\end{itemchoice}
	}
\end{ex}
\subsubsection*{PHẦN 3. Thí sinh trả lời từ câu 1 đến câu 6.}
\setcounter{ex}{0}
\begin{ex}% Câu 1
	Đun nước trong thùng bằng một dây nung nhúng trong nước có công suất $1,2\ kW$. Sau $3$ phút nước nóng lên từ $80^\circ C$ đến $90^\circ C$. Sau đó người ta rút dây nung ra khỏi nước thì thấy cứ sau mỗi phút nước trong thùng nguội đi $1,5^\circ C$. Coi rằng nhiệt tỏa ra môi trường một cách đều đặn. Bỏ qua sự hấp thụ nhiệt của thùng. Biết rằng nhiệt dung riêng của nước là $c = 4200\ J.kg^{-1}.K^{-1}$. Khối lượng nước đựng trong thùng bằng bao nhiêu $kg$ (làm tròn kết quả đến chữ số hàng phần trăm)?
	\SA[oly]{$3,55$}
	\loigiai{
	\begin{itemize}
		\item Ta có:
		$\begin{cases}
			\mathcal{P}.\tau _1 &= m.c.\Delta t_1 + \mathcal{P}_{hp}.\tau _1\\.
			\mathcal{P}_{hp}.\tau _2 &= m.c\Delta t_2 \Rightarrow  \mathcal{P}_{hp} = \dfrac{m.c\Delta t_2}{\tau _2}
		\end{cases}
		\Rightarrow \mathcal{P}.\tau _1 = m.c.\Delta t_1 + \dfrac{m.c\Delta t_2}{\tau _2}.\tau _1\\
		\Rightarrow 1,2.10^3.180 = m.4200.10+\dfrac{m.4200.1,5}{60}.180 \Rightarrow m \approx 3,55\ kg$.
	\end{itemize}
	}
\end{ex}

\begin{ex}% Câu 2
	\immini[thm]{
		Một học sinh làm thí nghiệm đo nhiệt hóa hơi riêng của nước. Ấm đun nước học sinh sử dụng có công suất $1500\ W$. Từ các kết quả thí nghiệm, học sinh vẽ được đồ thị quan hệ giữa khối lượng nước trong ấm và thời gian của quá trình hoá hơi của nước như hình vẽ. Nhiệt hoá hơi riêng của nước mà học sinh này đo được có giá trị là $x .10^6\ J/kg$. Xem như hiệu suất quá trình đun là $100 \%$. Tìm $x$.}
	{\begin{tikzpicture}[draw=Mapcolor, very thick, scale = 0.6, line join=round, line cap=round, >=stealth,line width=1.5]
			\draw[->] (0,0)--(0,5) node[right] {$m\ \left(g\right)$};
			\draw[->] (0,0)--(5.5,0) node[right] {$t\ \left(s\right)$};
			\draw[fill = black] (0,0) circle(2pt) node [below left] {O};
			\draw (0,4) node [left] {$300$} -- (5,1);
			\draw[dashed,line width=1] (0,2.8) node [left] {$200$} -- (2,2.8) -- (2,0) node [below] {$156$};
			\draw[fill=black] (2,2.8) circle (2pt) node [above] {M};
	\end{tikzpicture}}
	\par \shortans[oly] {$2,34$}
	
	\loigiai{
		\begin{itemize}
			\item $L=\dfrac{Pt}{\Delta m}=\dfrac{1500.156}{0,3-0,2}=2,34.10^6\ J/kg$
		\end{itemize}
	}
\end{ex}

\begin{ex}% Câu 3
	
	Trong máy nước nóng năng lượng Mặt Trời, năng lượng từ Mặt Trời được thu thập bằng nước lưu thông qua các ống trong bộ thu trên mái nhà. Bức xạ Mặt Trời đi vào bộ thu qua lớp vỏ trong suốt và làm nóng nước trong ống. Giả sử hiệu suất của toàn bộ hệ thống là $20\% $ (nghĩa là $80\%$ năng lượng Mặt Trời không dùng để làm nóng hệ thống). Cho biết nhiệt dung riêng của nước $4200\ J/kg.K$, khối lượng riêng của nước là $1000\ kg/m^3$. Để tăng nhiệt độ của $200$ lít nước trong bể từ $20^{\circ}C$ đến $40^{\circ}C$ trong $1$ giờ khi cường độ ánh sáng Mặt Trời tới là $700\ W/m^2$ cần diện tích tiếp xúc trực tiếp với ánh sáng Mặt Trời là bao nhiêu $m^2$ (làm tròn kết quả đến chữ số hàng đơn vị)?
	\shortans[oly]{$33$}
	\loigiai{
		\begin{itemize}
			\item $Q = mc\Delta t \Leftrightarrow \mathcal{H.P}\tau = D.V.c\Delta t \Leftrightarrow 0,2.700.S = 1000.200.10^{-3}.4200.20 \Rightarrow S \approx 33,33\ m^2$.
		\end{itemize}
	}
\end{ex}

\begin{ex}% Câu 4
	Bọ cánh cứng Bombardier Châu Phi (Stenaptinus insignis) có thể phun ra một luồng hơi nước từ phần đầu qua bụng để phòng thủ. Cơ thể bọ cánh cứng có các khoang chứa hai loại hóa chất; khi bọ cánh cứng bị đe dọa, các hóa chất này được kết hợp trong một khoang, gọi là khoang phản ứng, tạo ra một hợp chất được làm nóng từ $20^\circ C$ đến $100^\circ C$ bởi nhiệt lượng sinh ra từ phản ứng. Áp suất cao được tạo ra cho phép hợp chất được phun ra với tốc độ lên đến $19\ m/s$ ($68\ km/h$), làm hoảng sợ các loại động vật ăn thịt. (Bọ cánh cứng có chiều dài $2\ cm$). Giả sử nhiệt dung riêng của các hóa chất giống với nhiệt dung riêng của nước là $4,19. 10^3\ J/kg.K$ và nhiệt độ ban đầu của các hóa chất là $20^\circ C$. Nhiệt lượng sinh ra trên một đơn vị khối lượng ($1\ kg$) từ phản ứng của hai loại hóa chất bằng bao nhiêu $kJ$ (làm tròn kết quả đến chữ số hàng đơn vị)?
	\shortans[oly]{$335$}
	\loigiai{
		\begin{itemize}
			\item Nhiệt lượng sinh ra từ phản ứng là tăng nhiệt độ và chuyển hóa động năng: $$Q = mc\Delta t + \dfrac{1}{2}mv^2 = 1.4,19.10^3.80 + \dfrac{1}{2}.1.19^2 = 335380,5\ J \approx 335\ kJ$$.
		\end{itemize}
	}
\end{ex}

\begin{ex}% Câu 5
	Một thùng đựng $20$ lít nước ở nhiệt độ $20^{\circ}C$. Cho khối lượng riêng của nước là $1000\ kg/m^3$ và nhiệt dung riêng của nước là $4186\ J/kg.K$. Dùng một thiết bị điện có công suất $2,5\ kW$ với hiệu suất $70\%$ để đun lượng nước trên lên tới $70^{\circ}C$ thì thời gian truyền nhiệt lượng cần thiết là bao nhiêu giây?
	\shortans[oly]{$2392$}
	\loigiai{
		\begin{itemize}
			\item Ta có: $ Q_{ci} = mc\Delta t \Leftrightarrow Q_{tp}.\mathcal{H} = D.V.c.\Delta t \Leftrightarrow \mathcal{P}.\tau .\mathcal{H} = D.V.c.\Delta t \\
			\Rightarrow
			\tau = \dfrac{D.V.c.\Delta t}{\mathcal{P.H}} = \dfrac{1000.20.10^{-3}.4186.50}{2,5.10^3.0,7} = 2392\ s$.
		\end{itemize}
	}
\end{ex}

\begin{ex}% Câu 6
	Trong một bình nhiệt lượng kế có chứa nước đá nhiệt độ $t_1 = -5 ^\circ\text{C}$. Người ta đổ vào bình một lượng nước có khối lượng $m = 0,5\ kg$ ở nhiệt độ $t_2 = 80^\circ\text{C}$. Sau khi cân bằng nhiệt thể tích của chất chứa trong bình là: $V = 1,2$ lít. Biết khối lượng riêng của nước và nước đá là: $D_n = 1000\ kg/m^3$ và $D_{\text{đ}} = 900\ kg/m^3$, nhiệt dung riêng của nước và nước đá là $4200\ J/kg.K$, $2100\ J/kg.K$. Nhiệt nóng chảy riêng của nước đá là $340000\ J/kg$. Bỏ qua sự trao đổi nhiệt với bình nhiệt lượng kế và môi trường xung quanh. Khối lượng của chất chứa trong bình bằng bao nhiêu $kg$ (làm tròn kết quả đến chữ số hàng phần trăm)?
	\shortans[oly]{$1,18$}
	\loigiai{
		\begin{itemize}
			\item Nhiệt lượng do nước tỏa ra khi giảm về $0^\circ C$: $Q_n = m.c_n(t_2-0) = 168000\ J$.
			\item  Nhiệt lượng nước đá thu vào để tăng lên $0^\circ C$: $Q = m_{\text{đ}}.c_{\text{đ}}.(0-t_1) = 10500m_{\text{đ}}\ J$.
			\item Nhiệt thu vào để nước đá nóng chảy: $Q_{nc} = m_{\text{đ}}.\lambda = 340000m_{\text{đ}}\ J$.
			\item Nếu nước đá tan hết thì $m_{\text{đ}} = V.D_n - m = 1,2.1 - 0,5 = 0,7\ kg$.
			\item Với $m_{\text{đ}} = 0,7$ thì $Q<Q_n<Q_{nc} \Rightarrow $ nước đá không tan hết. Trong bình có cả nước và nước đá nên nhiệt độ cân bằng là $0^\circ C$.
			\item Theo đề bài, ta có:
			$\begin{cases}
				Q + Q_{nc} &= Q_n\\
				V_{\text{đ}} + V_n &=V 
			\end{cases} \Rightarrow 
			\begin{cases}
				10500m_{\text{đ}} + 340000m_{\text{đtan}} &= 168000\\
				\dfrac{m_{\text{đ}} - m_{\text{đtan}}}{D_{\text{đ}}} + \dfrac{m + m_{\text{đtan}}}{D_n} & = 1,2.10^{-3}
			\end{cases}\\
			\Rightarrow 
			\begin{cases}
				10500m_{\text{đ}} + 340000m_{\text{đtan}} &= 168000\\
				\dfrac{m_{\text{đ}}}{900} + \left( \dfrac{1}{1000} - \dfrac{1}{900} \right)m_{\text{đtan}}  & = 0,7.10^{-3}
			\end{cases} \Rightarrow
			\begin{cases}
				m_{\text{đ}} & \approx 0,677\\
				m_{\text{đtan}} & \approx 0,473
			\end{cases}
			$.
		\item Khối lượng chất trong bình là: $0,677 + 0,5 = 1,117 \approx 1,18\ kg$.
		\end{itemize}
	}
\end{ex}
